\section{Space Aliens, Pigs, Roosters, and Snakes}

\begin{quotex}
For we know that the whole creation groaneth and travaileth in pain together until now.
\flright{\textsc{Romans 8:22}}
\end{quotex}


Every now and then, this question comes up among some smart people: are there other planets in the Cosmos on which ``man'' did not fall? Obviously, this naively assumes that the notions of man and the universe are perfectly clear. It also assumes that a pre-fallen world is in historical, geographical, and scientific continuity its fallen counterpart.

On the contrary, the Fall of Adam affects all of creation, not just planet Earth. That seems to require that the pre-fallen world transcends the world as we now know it. Since no one believes in a flat Earth, then it is just as absurd to believe in a flat Cosmos. Rather, the esoteric teaching is that there is a vertical, as well as a horizontal, dimension to the Cosmos. Although the notion is common in alchemical and Hermetic texts, this is best illustrated in \textbf{Dante}'s \emph{Divine Comedy}.

\paragraph{The Three Poisons}

I used to spend many hours meditating on the Buddhist wheel of life, which represents all of creation. At its center, there are three animals representing the three poisons of the life of all sentient beings, not just earthbound humans. These are:

\begin{itemize}


\item \textbf{Ignorance}, represented by the pig


\item \textbf{Malice}, represented by the snake


\item \textbf{Concupiscence}, represented by the rooster

\end{itemize}


There are easily recognizable as the effects of the Fall of Adam. The Buddha recognized these poisons at the core of creation, but did not, and could not have, understood the original cause. In that scheme, the only answer is that they are the results of karma from previous lives. But does that not require an initial life at some previous point, a life that did not suffer from such defects?

Yet, even when he knows, and seems to understand, these poisons and their dangers, the chela finds he is still unable to overcome them. That is due to ``weakness'', the fourth effect of the Fall, but omitted in the Bhavacakra. The chela's guru will again attribute his weakness to a past life. Hence, the chela is stuck on the wheel.

Buddhism goes quite far in its understanding, but it falls short without the notion of grace, which is the antidote to weakness. Like the Buddha, Christ shows the way, but unlike the Buddha, he provides the Power.

As a practical matter, it has proven very difficult to even conceive of the life of a pre-fallen being. In popular culture, including the many movies and TV shows about aliens, they always act as though they had already consumed the poisons.

Moreover, the same applies a fortiori to androids. With all the chatter in vogue today about Artificial Intelligence threatening the human race, one would expect a better idea of its capabilities. An android, at best, would be capable of simulating the rational part of the human mind, although it would lack all interiority. Nevertheless, in films like \emph{Ex Machina}\footnote{\url{https://gornahoor.net/?p=8390}} and in HBO's \emph{Westworld}, the androids display all the effects of the poisons.

Both of those are examples of the fear engendered by the Sorcerer's Apprentice: man's creations will turn on him. It also shows how we judge by appearances, not by truth. ``If it looks human, it must be human.'' As I predicted some time ago (an easy prediction, of course), the first uses of androids will be for sex and for war. Sexbots are already for sale, and videos of dolls resembling attractive women are readily found. However, what will happen when sexbots are designed to simulate those few sexual acts that are still considered perverse in contemporary culture? Please, don't try to imagine it.

The danger is subconsciously recognized, although it does not arise from the androids, but rather from concupiscence not regulated by reality.

\paragraph{The Human Kingdom}


The Traditional view is that there are four kingdoms of manifest beings: Mineral, Plant, Animal, and Human. Flatland science regards human beings as just another species of animal life. That opinion is forced due to its limitation to appearances, DNA, homologous characteristics, in short, any and all exterior features.

When we turn to the qualitative features, i.e., interiority, the situation is quite different. Then, we see that the human being is not a type of animal, but a kingdom unto itself. A fortiori, any alien visitors from outer space, assuming they have these features, will be species of the human kingdom, no matter how their external form might appear.

Specifically, they will be fallen humans, subject to the three laws. The Eastern teachings from Yoga, the Vedanta, and the Tantras provide a detailed map of the inner life. Chart \ref{tab:MapInnerLife} shows the main features of that map, including the Western equivalent. I've slightly slanted it to be more aligned with Hermetic teachings, although without distorting it.

\begin{table}[h]\centering\small
\begin{tabular}{p{.2\textwidth}p{.35\textwidth}c}\toprule
\textbf{Yogic body} & \textbf{Sheaths} & \textbf{Western} \\\midrule Karana-sarira or causal body & Anandamaya kosha & Spirit \\\midrule Sukshma-sarira or subtle body &
\begin{tabular}{@{}p{.15\textwidth}p{.2\textwidth}@{}}\toprule
Vijnanamaya kosha & \begin{minipage}[t]{\linewidth}
\textbf{Intellect}\\
Intellective soul\\
Intuition\end{minipage} \\ 
Manomaya kosha & \begin{minipage}[t]{\linewidth}
\textbf{Mind}\\
Sensitive soul\\
Reason, ego, common sense\end{minipage}\\ 
Pranamaya kosha & \begin{minipage}[t]{\linewidth}
\textbf{Elan vital}\\
Vegetative soul \end{minipage}\\\bottomrule
\end{tabular} & Soul \\\midrule
Shula-sarira or gross body & Annamaya kosha & Body \\\bottomrule
\end{tabular}
\caption{Map of Inner Life}
\label{tab:MapInnerLife}
\end{table}

The three bodies correspond to three degrees of manifestation: gross, subtle, and formless. They can be regarded as sheaths over an inner core. The subtle body itself is further subdivided. Note that plant life has only the pranamaya sheath and animal life adds the manomaya kosha. The latter can be further divided into emotional and mental parts. The lower the life form, the less there is of the mental part.

The \textit{vihnanamaya} sheath, or intellective soul, is less easily recognizable without proper training. It is the center that grasps essences. The rational center, which most people assume to be the limit of intelligence, can only understand individual events, which it then constructs into a whole. Hence, in this view, a table is ``made'' of atoms, an animal is made of genes, and so one.

The intuitive power is synthetic and grasps the whole. In its view, the higher creates the lower, not vice versa. A human being differs radically from an android, because the human has an inner unity while an android is composed of parts. Its unity is built into from the outside.

So, an alien that is a metaphysical human being, would be composed of all those sheaths. Now a science fiction writer, who understands metaphysics, could do literary experiments with the bodies and sheaths. A very intelligent animal could conceivably lack the causal body and intellective soul, while having a highly developed rational part. It could easily pass the ``Turing test'', and there probably have been such beings on the earth. Attempts to create characters without an emotional center have not been very convincing.

\paragraph{Incarnation}


The mathematician and metaphysician, \textbf{Wolfgang Smith}, has done much to recover the traditional teaching of a multi-leveled cosmos, while integrating the concept with the latest scientific theories. For example, he shows that the body must be conceived not just as a corporeal entity, that is, as the shula-sarira in isolation. A lived body is energized by the vegetative soul, and that only is alive. Otherwise, the body is no more than a corpse. That soul is the body as experienced interiorly.

In addition, there is the physical substratum or body, i.e., the one studied by science solely through its quantitative techniques. The physical body is understood as neurons, biochemicals, and so on. Note that this is definitely not to be identified with the corporeal body which is in our daily experience. The pranamaya bonds with the corporeal, not the physical, body.

Now the Western teaching is that the intellectual soul (which would include the lower souls) bonds with the zygote at conception. The homolog to that in the pagan world was the story of Er who chose which body to incarnate into. The result of that compound is that the intellectual soul does not have full control over the lower forces, since the body has a separate inception. There is then a privation, since existence in the physical world resists the activity of the human essence. The human task is to bring more of his being under control of the Self. In psychological terms, the Self seeks to synthesize the unconscious with the conscious.

Note that Adam did not come from a zygote, so that his intellective soul had complete power over the body. This subtle body has four qualities: impassibility, agility, subtlety, and clarity\footnote{\url{http://www.newadvent.org/summa/4045.htm}}. This body can move at will by the power of thought, physical objects are impervious to it, and it does not suffer. You can get a sense of the subtle body in your dream life, if you pay attention. In a multi-leveled cosmos, the subtle body is a corporeal entity that exists apart from a physical substratum.

There is a lingering feeling, or longing, for that body as our true Being. There are many stories of saints who could levitate, bilocate, or read souls. The Chinese wuxia tradition is replete with Taoist masters with remarkable powers. Those are gained through an inner discipline of the mind and spirit. The \emph{Bene Gesserit} group in the Dune stories is not much different.

\paragraph{Alien Teachings}


So what can spacemen teach us earthlings? The schema just described does not depend on remarkable intelligence or high technical skill. These little green men will have the same problems dealing with the Fall that we have. There are two possibilities.

There are common sense tips, but we really don't need aliens for them. I have a birth certificate from Earth, and I can mention some. Don't get hooked on opioids, don't bomb children. That is merely a start. You can give the knowledge, but there is no easy cure for weakness.

Then, they could possibly provide us with extraordinary knowledge about philosophy, physics, black holes, disease prevention, economic systems, how to cook scrambled eggs (cream or no cream). But, seriously, how many people will understand? It is a futile endeavor, which I learned quite a while ago when I tried to teach freshman algebra and chemistry to education and phys ed majors.

There are two belief systems. The modern view is that all will be known in the indefinite future. The traditional view is that all important knowledge has been given to us already, but has been forgotten.

Which spiritual teacher will give us the answer to that dilemma?

Reference for multiple bodies: \textit{Neurons and Mind}, \textbf{Wolfgang Smith}

``Three Poisons'' image by Redtigerxyz -- \url{https://en.wikipedia.org/wiki/File:Bhavacakra_Thikse.jpg}

\flrightit{Posted on 2018-04-26 by Cologero}

\begin{center}* * *\end{center}

\begin{footnotesize}
\begin{sffamily}
\texttt{Lionel Chan on 2018-04-27 at 22:54 said:}

Can't wait for Wolfgang Smith's documentary. He should get together with the appropriately metaphysically primed parts of the Thunderbolts Project, would shake things up quite a bit as conventional cosmology continues to run into anomalies.

\hfill

\texttt{Byron Nightjoy on 2018-04-30 at 18:17 said:}

``Buddhism goes quite far in its understanding, but it falls short without the notion of grace, which is the antidote to weakness. Like the Buddha, Christ shows the way, but unlike the Buddha, he provides the Power''.

This is not true. In the teachings of the Buddhist sage, Shinran (1173-1263), one clearly finds a profound doctrine of grace based on the power of the Buddha (`tariki') which is given to us, as creatures of `blind passion', for our emancipation from samsara. In many ways, his insights take the notion of grace to its logical conclusion.

\hfill

\texttt{Cologero on 2018-04-30 at 18:33 said:}

Doesn't that just prove the point, Byron? Shinran taught Pure Land Buddhism and claimed that the traditional Buddhist practices are no longer effective. Hence, tariki is necessary to overcome weakness. In other words, one is no longer dominated by karma. Perhaps the traditions converge.

\hfill

\texttt{Byron Nightjoy on 2018-04-30 at 18:33 said:}

PS: An excellent introduction to this school of Buddhist thought (from a broadly traditionalist perspective) is \textit{Call of the Infinite: The Way of Shin Buddhism}, written by a priest of that tradition.

\hfill

\texttt{Byron Nightjoy on 2018-04-30 at 18:41 said:}

Hello Calogero. Yes, you're quite right -- `tariki' is quite necessary to overcome human weakness in light of `self-power' practices being no longer effective in this Dharma-ending age. However, unless we are able to surrender to the Buddha's grace, our lives will continue to be very much dominated by unfavorable karma and the realization of Nirvana thus becomes impossible for ordinary beings.

\hfill

\end{sffamily}
\end{footnotesize}

